\chapter{The MEDIA and DISCLOSURE}\label{ch:media}
\markboth{The Media and Disclosure}{The Media and Disclosure}

\chapquote{``The media's the most powerful entity on earth. They have the power to make the innocent guilty and to make the guilty innocent, and that's power. Because they control the minds of the masses.''}{Malcolm X}

This is where the book closes the loop.

We opened by dedicating this text to Arthur William Bell, and we opened Chapter~\ref{ch:language} by
arguing that the central problem in Ufology is not a shortage of sightings but a failure of language
--- that we cannot investigate a subject we cannot define. This final chapter is about the machine that
did the defining. Not the Air Force, not the scientists, not even the witnesses. The media.

Every term in this book reached you through a broadcast, a film, a newspaper column or a podcast.
``Flying saucer'' was a reporter's phrase. ``Little green men'' came from pulp fiction. The image in
your head when you read the word ``alien'' --- grey skin, black eyes, oversized cranium --- is a
specific visual design that spread through television and film in the late 1970s and 1980s, and
Section~\ref{sec:closeencounters}'s discussion of abduction narratives showed what happens to testimony once a template is
in circulation. \hi{The subject we have spent twenty chapters examining is, to a degree that should
bother us, a media artefact.}

So the last thing a sceptic needs is not another case file. It is an understanding of the delivery
system.

Before we take the machine apart, I want to introduce someone who has spent thirty years working
inside it on purpose, because he is the cleanest available demonstration of both halves of this
chapter's argument. The media distorts. The media is also, at present, the only instrument in
Ufology with a delivery mechanism that reaches the public at scale. Somebody has to stand in that
contradiction and work, and the person who has done it longest and most carefully is a documentary
film-maker rather than a scientist. That fact is itself a finding about the state of the field.

\figwrap[l]{0.36}{James Fox, whose seven documentaries have carried more first-hand UFO testimony to
general audiences than any peer-reviewed venue has managed in seventy-five years.}{https://commons.wikimedia.org/wiki/File:James_C._Fox.jpg}{fig:jamesfox}
James Fox was born on 25~April 1968 and came to the subject the way most people do --- through a
story told at home. His father had been aboard an aircraft in the 1950s whose crew watched something
they could not account for, and the account was told without embroidery and never monetised, which
is a combination Fox has spent his career trying to reproduce on film. He is not a scientist, holds
no relevant doctorate, and has never claimed one. \hi{His method is not analysis. It is
transcription.} He finds people who were present, puts a camera in front of them, and gets them to
say it once, on the record, with their names attached and their careers exposed. Chapter~\ref{ch:language}
would call that the correct first move: establish what was actually said before anybody argues about
what it means.

The filmography runs to seven features and it is worth listing in order, because the arc of it maps
almost exactly onto the changes in official language that this chapter ends with.
\emph{UFOs: 50~Years of Denial?} arrived in 1997 --- the fiftieth anniversary of Roswell, and a title
with a question mark that most of Fox's imitators would have dropped. \emph{Out of the Blue}
followed in 2003, premiering on the Sci~Fi Channel on 24~June 2003 with narration by the actor Peter
Coyote. \emph{I~Know What I~Saw} came in 2009 and was in effect \emph{Out of the Blue} rebuilt
around military and government witnesses alone. \emph{Pretty Slick} (2016) is not about UFOs at all
--- it is about the chemical dispersants sprayed after the Deepwater Horizon spill --- and I mention
it because it tells you what kind of film-maker he is when the subject cannot possibly be dismissed
as fantasy. \emph{The Phenomenon} (2020), produced with Marc Barasch, brought Jacques
Vall\'ee, the Ariel School children, Rendlesham and Senator Harry Reid into a single
narrative. \emph{Moment of Contact} (2022) is a full-length treatment of Varginha. \emph{The
Program} (2024) turns the camera on the alleged retrieval and reverse-engineering programmes that
Section~\ref{sec:immaculate} deals with in documentary terms.

\emph{Out of the Blue} is the one to understand, and I will be blunt about why. It is the film that
established the template every serious UFO documentary since has either followed or reacted against:
no reconstructions of alien surgery, no theremin on the soundtrack, no narrator asserting what the
objects were. Fox went instead to the people with the most to lose --- Air Force missile-launch
officers, air traffic controllers, a former Chief of Staff of the Belgian Air Force, astronauts,
Gordon Cooper, the Iranian pilot from the 1976 Tehran incident --- and let them describe their own
instrument readings. The Minuteman shutdown material in Chapter~\ref{ch:nuclear} is on film largely
because Fox and his collaborators went and asked. \hi{The film's argument is procedural rather than
extraterrestrial: credentialed people filed reports through official channels, and the channels did
nothing with them.} That is this book's thesis arriving from the other direction, in 2003, on cable
television, fourteen years before the \emph{New York Times} made it respectable.

\begin{funfact}
\emph{Out of the Blue} is routinely dated to 2002 in citations, filmographies and even on some
official-looking pages. The theatrical and festival material dates from 2002, but the release the
world actually saw --- the Sci~Fi Channel broadcast that made the film's reputation --- went out on
24~June 2003. I have made this error myself in print. It is a small demonstration of
Section~\ref{sec:warworlds}'s lesson: a date enters circulation, gets copied rather than checked, and
survives on repetition alone.
\end{funfact}

Now the sceptical accounting, because the point of putting Fox at the front of this chapter is not to
praise him. A documentary is an argument with a budget, and it has the three weaknesses that come
with the form. It selects: Fox chooses which witnesses appear, and the ones who make weak or
confused claims are edited down or omitted, which produces a corpus of testimony that looks more
consistent than the raw record is. It scores: music, cutting rhythm and the ordering of interviews
do persuasive work that no footnote can be attached to. And it cannot cross-examine --- a witness on
film is never asked the hostile follow-up question, because the film needs the witness to keep
talking. \hi{Fox's material is testimony, not data.} It is very good testimony, gathered with more
care than most of what this field produces, and it still sits on the lower rungs of the evidential
ladder that Chapter~\ref{ch:language} set up.

\figwrap[r]{0.40}{The memorial to the ``ET de Varginha'' in Varginha, Minas Gerais. The case Fox
filmed in \emph{Moment of Contact} (2022) has become municipal identity: the town now has a
spacecraft-shaped water tower and a memorial to the entity three young women said they saw. No
photograph of the 1996 object or being exists; this is what a case looks like after it has been
absorbed by the culture that hosts it.}{https://commons.wikimedia.org/wiki/File:Memorial_do_ET_de_Varginha.._19_de_fev._de_2025.jpg}{fig:varginhamem}
\emph{Moment of Contact} illustrates both edges of the blade at once. Varginha, which
Chapter~\ref{ch:intro} treats as the one case I cannot file away tidily, is a case built almost
entirely on human statement --- the three young women, the military drivers, the doctors, the
six-hundred-page 2010 Army inquiry. Fox got a great many of those people on camera, several of them
for the first time, and that is a genuine contribution to the record which would not otherwise
exist. He also filmed them in a town that has spent thirty years constructing its own economy around
the story, complete with the memorial in Figure~\ref{fig:varginhamem}, and there is no honest way to
measure how much of the 2022 testimony is 1996 memory and how much is thirty years of local
retelling. Both things are true and the film can only really hold one of them.

The same double reading applies to \emph{The Phenomenon} and the Ariel School sequence.
Chapter~\ref{ch:intro} notes that Cynthia Hind did the one thing that matters with child witnesses
--- separating them and having them draw before they could compare notes --- and that convergent
child testimony is simultaneously powerful and the easiest thing in the world to contaminate. Fox
returned to those witnesses as adults, decades later, and filmed them. That is valuable. It is also,
unavoidably, a fourth or fifth pass over memories that have been publicly rehearsed since 1994.

So why open the chapter about media distortion with a film-maker? Because Fox is the working answer
to the question this chapter poses. The distortion documented in Sections~\ref{sec:warworlds}
through~\ref{sec:influencehollywood} is not an argument for abandoning media; it is an argument for
using it deliberately. \hi{Fox's contribution is not proof of anything in the sky. It is the
demonstration that a first-hand account can be moved from a witness to a mass audience without being
turned into entertainment on the way ---} which, after twenty chapters on attenuation, is close to
the most hopeful sentence in this book. He also, and this is the part his critics never mention,
changed what he was willing to say as the record changed rather than the other way round, which is
the behaviour Section~\ref{sec:doubtdrivenskepticism} called doubt-driven scepticism and which is rarer in
this field than any sighting.

\pullquote[James Fox, on his method across seven films]{``I don't ask people what they believe. I ask
them what they saw, and then I let the tape run.''}

That instinct --- let the tape run, let the witness finish --- is the same instinct we will meet in
Section~\ref{sec:artbell} in an entirely different medium, and it is the note this book ends on. Keep
it in mind for the next sixty pages, because almost everything else in this chapter is a catalogue of
what happens when nobody bothers.

% ---------------------------------------------------------------
\section{War of The Worlds}\label{sec:warworlds}

On the evening of 30 October 1938, the Columbia Broadcasting System aired an adaptation of H.\,G.\
Wells's \emph{The War of the Worlds}, directed by and starring the twenty-three-year-old Orson Welles.
The innovation was structural: rather than narrating a story, the first two-thirds of the broadcast
imitated the form of a live radio evening --- dance music from a hotel ballroom, interrupted by
increasingly urgent news bulletins, remote reports from Grover's Mill, New Jersey, and an on-air
transmission that cut to dead silence.

What happened next is the most repeated media anecdote of the twentieth century, and most of what you
have heard about it is wrong. This matters more than the story itself, so let us do it properly.

\figwrapfix[l]{0.38}{Orson Welles at the CBS microphone in 1938, headphones on and pipe in hand, during the period of the \emph{Mercury Theatre} broadcasts. The photograph shows the studio, not the panic --- the panic was manufactured the following morning, in print.}{https://commons.wikimedia.org/wiki/File:Orson_Welles_1938.jpg}{fig:welles}
The legend says a nationwide panic ensued (see Figure~\ref{fig:welles}): millions fleeing, highways jammed, suicides, mobs.
Subsequent scholarship --- particularly Jefferson Pooley and Michael Socolow's work, and the media
historian W.\ Joseph Campbell's --- has established that the audience was small. The broadcast ran
against \emph{The Chase and Sanborn Hour} with Edgar Bergen, which was the most popular programme in
America, and the C.\,E.\ Hooper ratings service found roughly two per cent of surveyed households
listening to Welles. There was no confirmed death, no verified suicide, and the ``jammed highways''
appear in almost no contemporaneous local reporting. Some listeners were genuinely frightened, some
called their local stations, and that is approximately the extent of it.

So where did the panic come from? Overwhelmingly from newspapers. The next morning's press ran the
story hard and kept running it, and there is a straightforward commercial explanation: radio was
five years into eviscerating newspaper advertising revenue, and a story establishing that the new
medium was dangerous and irresponsible was extremely useful to the old one. The panic was not an
event that the press covered. It was a product the press manufactured, out of an incentive it had.

\begin{funfact}
There is a delicious recursion here. The single most cited example of ``the media can make people
believe anything'' is itself a thing the media made people believe. If you want one case study to
keep in your pocket for the rest of your life, it is this one: a story that is famous for
demonstrating credulity, which almost everyone who repeats it has accepted uncritically, and which
survives because it is too good a parable to check.
\end{funfact}

The genuinely important consequence is that the broadcast established a \emph{format}. Welles proved
that the conventions of authoritative news --- the calm anchor, the remote correspondent, the
interruption --- carry credibility independent of content. Nine years later, when the phrase ``flying
saucer'' entered the wire services after Kenneth Arnold, the machinery for turning an ambiguous
sighting into a national phenomenon overnight was already built and tested.

\begin{srcnotes}
\item Pooley, J.\ \& Socolow, M.\ (2013). The Myth of the War of the Worlds Panic. \emph{Slate}.
\item Campbell, W.\ J.\ (2010). \emph{Getting It Wrong: Ten of the Greatest Misreported Stories in
American Journalism}. University of California Press.
\end{srcnotes}

% ---------------------------------------------------------------
\section{Power of Propaganda}\label{sec:powerpropaganda}

I want to use this section narrowly and technically, because ``propaganda'' has become a word people
throw at coverage they dislike. What I mean by it is specific: the deliberate shaping of belief
through control of framing and repetition rather than through falsehood. The most effective
propaganda contains very few lies. It works by selection.

Four mechanisms are worth naming, because you will find all four in the UFO literature and in the
official response to it.

\textbf{Framing.} The same facts produce different conclusions depending on the container. ``Air Force
releases report on Roswell debris'' and ``Air Force issues second explanation for Roswell'' describe
an identical event. Neither is false. Chapter~\ref{ch:roswell} is a case study in framing: the 1994 and
1997 reports were reasonable documents that arrived so late, and after so much prior official
vagueness, that their framing was decided before they were read.

\textbf{The illusory truth effect.} Repeated statements are rated more likely to be true, and this
occurs independently of the source's credibility and independently of the listener's prior knowledge.
This is well replicated experimental psychology, not commentary, and it is the single most important
finding in this chapter, because it means that the sheer \emph{volume} of UFO content generates
belief mechanically. No one has to lie. Something merely has to be said often.

\textbf{Availability and the vividness advantage.} We estimate frequency by how easily examples come to
mind. A dramatic televised abduction account is retrieved effortlessly; a statistical base rate is not
retrieved at all. \hi{This is why one good story outperforms a hundred good studies.}

\textbf{Selective silence.} Section~\ref{sec:aviationsecrecy} covered what military secrecy actually
requires: a programme must not merely be hidden, its absence must be explained. When an agency refuses
to comment, it is not producing an absence of information --- it is producing information, and the
audience will fill the gap with the most interesting available hypothesis. The Air Force's handling of
the U-2 and the Blackbird did more to build the flying-saucer mythology than any believer, and it did
it by saying nothing at all. Secrecy is not the opposite of publicity. It is a form of it.

\pullquote{``Propaganda rarely needs to tell you something false. It only needs to decide which true
things you hear, and how often.''}

And because I have been asking you all book to apply these tools to believers, apply them here to me.
This book selects. I chose which cases to include, I put the debunked crystal skulls in a full section
and compressed six decades of sonar contacts into a few paragraphs, and I framed each chapter with a
quotation I picked. I have tried to flag my uncertainty where it exists, but the structure of a
textbook is itself a framing device, and you should read this one the way I have taught you to read
everything else in the field.

% ---------------------------------------------------------------
\section{The Influence of Hollywood}\label{sec:influencehollywood}

I should declare an interest before I argue anything in this section, because it would be dishonest
not to. I did not arrive at this subject through the archives. I arrived through screens. Superheroes,
cartoons and science fiction are the frame I have used my whole life to make sense of the world ---
they are how I first learned what a hero looks like, what a secret looks like, what an alien looks
like, and what a government that knows something looks like. Every reader of this book has a
frame like that whether they admit to it or not, and mine is Saturday morning television and a
rented cinema seat.

My hero is Spider-Man, and I do not think that is incidental to the rest of this book. Peter Parker
is not a god, an alien or a billionaire. He is a teenager who gets bitten by something in a lab, is
handed evidence that the world does not work the way he was taught, and has to decide what to do
about it while nobody in authority helps him. He hides what he knows, and the hiding costs him
everything he loves. \hi{The line every child remembers --- with great power comes great
responsibility --- is, read cold, a statement about disclosure.} It says that possessing knowledge
other people do not have creates an obligation, not a privilege. That is the exact ethical claim
Chapter~\ref{ch:government} makes against the classification system, and I first met it in a comic.

It is worth noticing that almost every hero I grew up with is an atomic-age child: radiation,
irradiated bites, gamma bombs, unstable isotopes, accidents in government laboratories. The same
anxiety that Chapter~\ref{ch:nuclear} traces through the sighting record --- that we had built
something we did not understand and that men in uniform were keeping the details from us --- is the
engine of the entire genre. The comics and the flying saucers are not two phenomena. They are one
culture metabolising one fear in two formats.

So when I say the traffic between screen and sighting runs in both directions, I am not standing
outside it. I am a data point in it. Now, the chronology, which is the argument.

The pattern begins immediately. \emph{The Day the Earth Stood Still} and \emph{The Thing from Another
World} both appeared in 1951, four years after Arnold, and between them established the two enduring
templates: the benevolent emissary with a warning, and the infiltrating menace. Cold War anxiety wore
a flying saucer for a decade after that --- \emph{Invasion of the Body Snatchers} in 1956 is a film
about aliens in the same sense that \emph{Moby-Dick} is a book about whaling.

Then come the two that changed the field. \emph{Close Encounters of the Third Kind} in 1977 took its
title and its classification scheme directly from J.\ Allen Hynek, who consulted on the film and
appears in it, and it did something no previous film had: it made contact benign, awe-inspiring, and
quasi-religious. Chapter~\ref{ch:contact}'s discussion of what contact would actually involve is in
large part a conversation with this film, because Spielberg's version --- lights, music,
communication as a five-note motif --- is the version in most people's imagination.

\emph{Fire in the Sky} in 1993 dramatised the Travis Walton case, and \emph{The X-Files}, running from
1993 to 2002, did something more consequential than any single film. It made ``the government is
hiding it'' the default frame rather than a fringe position, gave the culture the phrase ``trust no
one'', and coincided almost exactly with the peak years of \emph{Coast to Coast AM}. The most
influential UFO text of the 1990s was a network drama, and the second most influential was a radio
show.

\ufofloatnext
\figwrap[r]{0.38}{The 1987 cover art of \emph{Communion} --- the point after which reported entities converge on a single face.}{https://commons.wikimedia.org/wiki/File:Communion_book_cover_art.jpg}{fig:communion}
Here is the part that should trouble a serious investigator. \emph{Communion}, Whitley Strieber's
account, was published in 1987 and its cover carried a specific rendering (see Figure~\ref{fig:communion}) of a grey-skinned,
large-eyed being. Abduction reports from before that period describe a much wider range of entities:
hairy dwarfs in the South American reports, tall Nordic figures, robotic forms, one-off oddities. The
reports converge on the \term{Grey} \emph{after} the image saturates the culture, and they converge
further after \emph{The X-Files}. Section~\ref{sec:hypnosis} laid out the mechanism --- suggestibility under
hypnosis, source-monitoring error, the reconstructive nature of memory --- and this is the
epidemiological evidence for it. Testimony tracks the movies.

\begin{funfact}
Before \emph{The Day the Earth Stood Still}, ``the aliens'' in American popular culture were
frequently invertebrate, gaseous, or unrepresentable. The reason your mental image is humanoid,
bipedal, and roughly your own height is production economics: it is far cheaper to put an actor in a
suit. The dominant visual hypothesis about extraterrestrial biology in the twentieth century was
determined by studio budgets.
\end{funfact}

None of this proves that no one saw anything. It proves that the \emph{description} of what they saw
is not independent evidence, and independence is the entire value of corroboration. \hi{A thousand
witnesses who all watched the same film are, evidentially, closer to one witness than to a thousand.}

\subsubsection{Stargate: the base that is real and the door that is not}

If one property has done more than any other to shape how I --- and, I would argue, an entire
generation --- picture the machinery of official secrecy, it is \emph{Stargate}. Roland Emmerich and
Dean Devlin's film arrived in 1994 with a premise lifted almost intact from Chapter~\ref{ch:ancient}:
a ring of unknown metal is dug up at Giza in 1928, it turns out to be a transport device, and the
Egyptian god Ra turns out to have been an extraterrestrial wearing a civilisation as a costume. That
is Erich von D\"aniken's thesis with a budget. The film did not invent the ancient-astronaut
hypothesis; it did something more effective, which was to make it look like an engineering problem
rather than a theory --- something you could point a camera at, staff with soldiers, and file reports
about.

Then \emph{Stargate SG-1} ran from 1997 to 2007 across ten seasons and 214 episodes, with
\emph{Atlantis} and \emph{Universe} after it, and it did the thing that matters for this chapter. It
put the door under Cheyenne Mountain.

Read Section~\ref{sec:cheyennemountain} again with the television in mind. The Cheyenne Mountain
Complex is real. It is 600 metres of granite, fifteen buildings on half-tonne springs, two 25-tonne
blast doors, its own water reservoir, its own power. Every one of those facts is published, and I
gave the construction dates. What \emph{SG-1} did was take that genuinely existing, genuinely
documented hole in a mountain and install a wormhole on level 28, and it did so convincingly enough
that for a very large number of people the phrase ``Cheyenne Mountain'' now carries a faint
radioactive glow of the extraterrestrial that nothing in the actual record put there.

\actualq[problem]{that a television show made up a secret base.}{that it did not have to. The base was
already there, already deep, already armoured, already run by men who genuinely could not tell you
what was inside. Fiction only had to add one prop to a facility the public had been handed for free.}

\figwrap[r]{0.34}{A screen-used Horus guard costume and staff weapon from the \emph{Stargate} productions, on museum display in Lyon. The iconography is Egyptian; the premise is von D\"aniken's.}{https://commons.wikimedia.org/wiki/File:Stargate_prop.jpg}{fig:stargateprop}
This is the mechanism I want you to hold onto, because it is subtler than ``Hollywood invents
conspiracy theories''. Hollywood mostly does not invent. It \emph{annexes}. It finds a real
institution with a real secrecy gradient --- a mountain, an air base, a records office --- and it
fills the part the public cannot see with narrative, and because the visible half is verifiable the
invisible half inherits its credibility. Chapter~\ref{ch:historyarea51} is the same trick played on a
dry lake bed in Nevada. Section~\ref{sec:dulce} is the same trick played by amateurs on a mesa in New
Mexico. The load-bearing element is never the invention. It is the true part underneath it.

There is one more coincidence in the name that I have to address, because readers who reach
Chapter~\ref{ch:abduction} and then think about this section will notice it and assume I missed it.
The United States really did run a programme called \hi{Project Star Gate}. It was the remote-viewing
effort at Fort Meade --- the one that came out of Grill Flame, Center Lane and Sun Streak, that took
the Star Gate name around 1991, and that was closed down in 1995 after an external review found it
had produced nothing of actionable intelligence value. The codename was in use before Emmerich's film
reached cinemas, and I have found nothing whatever to connect the two beyond the words. It is a
collision, not a clue. But it is a beautifully instructive collision, because I can tell you from
experience what happens when you say ``the government had a project called Stargate'' out loud to
someone whose only reference is the show: they hear confirmation. \hi{Two unrelated things sharing a
name is one of the cheapest and most reliable ways for a false belief to enter the record}, and
Section~\ref{sec:hypnosis} and Chapter~\ref{ch:language} together should have taught you to hear the
alarm bell on it.

I want to be fair to the show as well as sceptical of its residue, because I love it. \emph{SG-1} is
the rare piece of screen ufology that models good practice: a team with a linguist and an
astrophysicist on it, mission reports, a debriefing room, a habit of being wrong in the first act and
revising by the third. Whatever it did to Cheyenne Mountain's reputation, it also spent ten years
showing a mass audience that investigating the unknown is paperwork, argument and expertise rather
than a lone believer with a flashlight. I would rather a generation absorbed that than the
alternative.

\begin{funfact}
The \emph{Stargate} franchise ran for seventeen consecutive years across three live-action series and
held a Guinness record for the longest-running consecutive science-fiction television series in North
America, at ten seasons, before that mark moved on. For most of its life it was broadcast on a channel
that also carried programmes presenting Cheyenne Mountain and Area~51 as documentary subjects, in the
same schedule, with the same title cards. This chapter's whole argument about the collapse of the
label is visible in a single evening's listings.
\end{funfact}

\subsubsection{Cartoon Aliens}

Animation deserves its own heading here, and not as a footnote to film, because animation is where
the alien got out of the rubber suit.

\emph{Ben 10} premiered on Cartoon Network at the end of December 2005, created by the four-man
collective Man of Action --- Duncan Rouleau, Joe Casey, Joe Kelly and Steven T.\ Seagle. The premise
is a ten-year-old boy, Ben Tennyson, on a summer road trip in his grandfather's camper van with his
cousin Gwen, who finds an alien device that fuses to his wrist. The Omnitrix lets him become ten
species at will, and the original ten are the point of this section: a being made of fire, a walking
crystal, a four-armed strongman, a blind orange quadruped that hunts by smell, a tiny grey amphibian
with a large brain, a velociraptor built for speed, an insectoid, an aquatic predator, a living
software patch, and a semi-transparent phantom.

\figwrap[l]{0.34}{A child wearing an Omnitrix toy --- the wristband that turns a ten-year-old into ten different non-humanoid species. No suit budget, no humanoid constraint.}{https://commons.wikimedia.org/wiki/File:Ben_10_watch.jpg}{fig:ben10}
Now set that list beside the fun fact earlier in this section. The reason twentieth-century screen
aliens are humanoid, bipedal and roughly our own height is that a costume is cheaper than a creature.
Animation does not pay that tax. \hi{The moment the suit budget vanishes, alien morphology explodes
--- and it explodes in exactly the directions astrobiology says it should}: silicon, plasma,
radically different sensoria, body plans with no bilateral symmetry, intelligence not sited in a head.
Chapter~\ref{ch:life} spends its length arguing that convergent evolution does not licence us to expect
something shaped like us. A children's cartoon from 2005 had already staged that argument, brightly,
twenty-six minutes at a time.

And then look at the reported record. The entities in abduction testimony did not diversify when
animation showed the culture how. They converged --- on the Grey, on the humanoid, on a body that
would fit an actor. The genre with unlimited creature freedom produced ten wildly different
non-humanoids for one show, while the eyewitness literature, which is supposedly reporting reality,
produced one face over and over. \hi{That asymmetry is evidence about the testimony, not about the
universe.}

\emph{Ben 10} is the primary case because it is the most systematic --- it is functionally a
bestiary in weekly instalments, and it kept going through \emph{Alien Force} (2008),
\emph{Ultimate Alien} (2010), \emph{Omniverse} (2012) and a 2016 reboot. But it sits in a long line.
Marvin the Martian arrived in 1948 and taught American children that the invader might be small,
pompous and funny. \emph{Invader Zim} (2001) made the advance scout an incompetent embarrassment to
his own empire. \emph{Lilo \& Stitch} (2002) recast the escaped experiment as family.
\emph{Futurama} (1999) and later \emph{Rick and Morty} (2013) turned the galactic bureaucracy into
satire, and \emph{Steven Universe} (2013) built an entire occupation narrative out of mineral
lifeforms. Every one of these is doing serious work on how a generation imagines contact, and none of
them is ever cited in the ufological literature, because cartoons are assumed not to count.

They count. The image in a witness's head when they look up did not come from a peer-reviewed paper.

\begin{funfact}[Fun Fact: The Author Declares His Bias]
\emph{Ben 10} was my favourite television show as a child. I am not going to pretend otherwise in a
book that spends twenty chapters demanding other people disclose their priors. A boy with a device on
his wrist that reveals hidden forms of life, travelling with a grandfather who turns out to have known
about aliens the whole time and not told him --- if you want to understand why the author of this
textbook finds the disclosure question emotionally as well as intellectually compelling, you may start
there.
\end{funfact}

One honest note about the illustration. \emph{Ben 10} is Cartoon Network property, and there is no
freely licensed still, character design or frame from the series that I could use here --- the
canonical images are all under copyright and the only Commons material that touches the franchise is
merchandise photographed by members of the public. So the figure above is a licensed photograph of a
child wearing an Omnitrix toy rather than anything from the show itself, which is, if anything, the
more relevant picture: it is a photograph of the transmission mechanism. That is how the aliens
actually got into a generation's imagination --- not through the archives, but through the toy aisle.

\subsubsection{Aliens on YouTube}

Hollywood needed a studio, a distributor and a release date. Cartoon Network needed a network. The
third era needs a phone. \hi{Once upload was free, the pipeline that had produced every screen alien
for eighty years was replaced by one with no budget, no credit roll and no obligation to tell you it
was fiction} --- and that last removal is the whole change, because a film announces itself as a film
and a YouTube clip announces itself as a record.

\textbf{The Polished Knob} is the primary case study for this section, and it is chosen deliberately
over the famous hoaxes. The channel posted a run of short handheld clips beginning around 2010 ---
night footage, a fixed vantage, a great deal of shouting, and in the best-known sequences a tall,
thin, pale humanoid figure moving at the edge of the light. Around 219 videos and roughly 28,000
subscribers is a small operation by platform standards. The reach is not.
One clip alone has been watched over 900,000 times, and the figure from it --- widely described in
circulation as a ten-foot Grey --- has since been cut, re-cut, reposted and narrated across TikTok and
Instagram by accounts that never name the source at all. Original material is available at
\srclink{https://www.youtube.com/@ThePolishedKnob}, and it should be watched, because reading about it
is not the same as noticing what it does to you.

What makes it the right case study is the reception, not the footage. Trace the discourse and you find
every pathology in this book compressed into a comment section. Fringe sites list the channel as ``a
catalogue of numerous genuine original sightings'' with no analysis whatsoever. Sceptics dismiss it as
an obvious hoax, also with no analysis whatsoever --- and one of the more honest threads on the
subject, posted to \emph{r/aliens} in 2023 under the title ``Reconsidering ThePolishedKnob'', makes
precisely that complaint: the material was rejected as fake with about as much evidence as was offered
for it being real. Meanwhile the clips themselves have drifted loose of their origin, so that footage
from unrelated comedy creators now gets reposted \emph{as} Polished Knob material and Polished Knob
footage gets reposted under other names. \hi{Nobody in the chain --- believer, debunker or reposter ---
has done the one piece of work that would settle it: examined the original file.} Fifteen years on,
there is still no published frame analysis, no lighting reconstruction, no attempt to locate the site,
no interview with the uploader on the record.

That is the finding. The platform did not lower the standard of evidence in ufology; it exposed that
there was never a standard, only a gatekeeper. When a television network stood between a claim and its
audience, somebody at least had to be persuaded before it aired --- badly, commercially, but somebody.
Remove the network and the claim reaches ten million phones with nothing in between, and the
field's response is to sort it by vibe. \hi{A generation's baseline image of an alien encounter is now
set by whichever uploader is most watchable, and watchability is not a truth criterion.}

The technical asymmetry has also inverted since the era described earlier in this chapter. A 1975
television alien cost a costume department a fortnight; a 2025 one costs a laptop and an afternoon,
and the current generation of generative video tools will produce a plausible night-vision humanoid
from a text prompt in seconds. The Polished Knob material predates that capability, which is exactly
why it is worth studying now: it is the last generation of footage where the ``it could be CGI''
dismissal actually cost something to be true, and the field \emph{still} did not do the work. Anything
uploaded after this point arrives into an environment where fabrication is free and verification is
not, and there is no reason to expect the analysis to improve when the volume goes up by four orders
of magnitude.

\begin{funfact}
Note which way the causation runs. The tall, thin, pale, long-limbed figure in this footage is not a
departure from the screen tradition described in this chapter --- it is the tradition, rendered by
someone who grew up inside it. Whitley Strieber's \emph{Communion} cover, the Greys of 1990s
television and two decades of film all arrive before the camera does. \hi{Whatever is or is not in
these videos, the shape people report seeing in them was issued to them by Hollywood first.}
\end{funfact}

% ---------------------------------------------------------------
\section{The Freedom Of Information Act}\label{sec:foia}

Every previous section of this chapter has been about people trying to get the public's attention.
This one is about the only tool ufology ever had for getting the \emph{government's} attention, and
it is not a camera, a witness or a documentary. It is a form letter.

The Freedom of Information Act was signed on 4 July 1966 and took effect a year later. Its premise
reversed three centuries of administrative habit: a citizen no longer had to demonstrate a need to
know before seeing a federal record. The agency had to demonstrate a legal reason to withhold it.
There are nine such reasons, the exemptions, and the two that matter to this book are \textbf{b(1)},
properly classified national-defence information, and \textbf{b(3)}, material another statute forbids
an agency to release. Everything else in the file is supposed to come out.

As written in 1966 the Act had no teeth. Agencies could take years to answer, charge whatever they
liked, and treat a refusal as final; judges were expected to defer to the executive's word that a
secret was a secret. Congress fixed this in the 1974 amendments, passed in the aftermath of Watergate
over President Ford's veto. They imposed deadlines, capped fees, allowed courts to award legal costs
against the government, and --- the decisive change --- authorised federal judges to examine the
withheld records themselves, \emph{in camera}, and rule on whether the classification was justified.
\hi{From 1975 onward a private citizen with a typewriter and a persistent lawyer could compel an
intelligence agency to explain itself to a judge.} That is the machinery on which the documentary
history of this subject was built.

\subsection{Ground Saucer Watch and the CIA's nine hundred pages}\label{sec:gsw}

The first great demonstration came from a small Phoenix organisation called Ground Saucer Watch,
which in 1977, represented by attorney Peter Gersten, sued the Central Intelligence Agency for its
UFO records. The CIA's public position for two decades had been that it held essentially nothing:
some interest around 1952, the Robertson Panel of January 1953 described in
Chapter~\ref{ch:government}, and no continuing programme.

\figwrap[r]{0.38}{A declassified CIA page released under the Freedom of Information Act. The information
is in the document; the document is in the reader's hands; and the reader still does not have the
information. This is what a win looks like.}{https://commons.wikimedia.org/wiki/File:Redacted_CIA_document.jpg}{fig:foiaredact}
In December 1978 the Agency released roughly nine hundred pages. \emph{The New York Times} reported
the outcome on 14 January 1979 --- and the story was not that aliens were confirmed. The story was
that the CIA had maintained an interest in the phenomenon long after it said it had stopped, had
collected reports from abroad, had run photographic and instrumentation analysis, and had told the
public otherwise. Fifty-seven documents were withheld outright on b(1) and b(3) grounds, and the
released pages were perforated with redactions of the kind shown in Figure~\ref{fig:foiaredact}.

That pattern --- release confirms concealment, concealment survives release --- is the single most
important lesson in the FOIA chapter of this subject, and it is why serious researchers stopped
asking agencies whether they had a UFO programme and started asking for the paperwork a programme
necessarily generates: routing slips, filing indices, teletype logs, travel vouchers. Bureaucracy is
more honest than testimony because bureaucracy is not trying to convince anyone of anything.

\subsection{CAUS versus the National Security Agency}\label{sec:causnsa}

Gersten's next client was Citizens Against UFO Secrecy, and his next defendant was harder. CAUS sued
the National Security Agency in 1980 for its UFO holdings. The NSA acknowledged 156 relevant
documents and refused 135 of them, supporting the refusal with an \emph{in camera} affidavit filed on
9~October 1980 by Eugene Yeates, the Agency's Chief of Policy. The affidavit explained to the judge,
and to nobody else, why the material could not be described, let alone released. The court read it
and ruled for the Agency. Under FOIA, an argument the public is not permitted to hear can still be
the argument that defeats the public.

A redacted version of the Yeates affidavit was finally released in 2005, twenty-five years later, and
vindicated the sceptics in the least satisfying way available: the withheld material was almost
entirely signals-intelligence traffic --- foreign communications the NSA had intercepted, in which
other countries discussed UFO sightings. The secret was not what was seen. The secret was that the
United States was reading the mail of the people discussing it. \hi{The classification protected the
method, not the phenomenon}, which is the answer in a majority of cases where a UFO file is
sealed --- and it is an answer that took a quarter-century and a federal lawsuit to obtain.

\begin{funfact}[Fun Fact: The Glomar Answer]
In 1975 a journalist asked the CIA about the \emph{Glomar Explorer}, a Howard Hughes vessel that was
secretly raising a sunken Soviet submarine. The Agency replied that it could ``neither confirm nor
deny'' the existence of the records requested, because to do either would itself reveal something.
The courts upheld it, and the \term{Glomar response} entered federal law as a legitimate reply to a
FOIA request. Every time an agency tells a UFO researcher it can neither confirm nor deny, it is
citing a case about a submarine, not a saucer --- and it is not being evasive in the colloquial
sense. It is invoking a precedent.
\end{funfact}

\Needspace*{0.40\textheight}
\subsection{What the Act could not do}\label{sec:foialimits}

\figwrap[l]{0.36}{A page of the FBI's Majestic~12 file, obtained under the Freedom of Information Act
and now published in the Bureau's online Vault. Everything about it is authentic --- the file, the
stamps, the release --- except the document inside it, which the FBI itself marked
``BOGUS.''}{https://commons.wikimedia.org/wiki/File:Majestic_12_Part_1_of_1_(page_4_crop).jpg}{fig:mj12memo}
FOIA is a records law, not a truth law, and the distinction ruins more arguments than any other single
fact in this book. A document released under the Act is proof that the government held it. It is not
proof that the contents are true, that the author was informed, or that anyone above the filing clerk
ever took it seriously. The Majestic~12 papers of Section~\ref{sec:mj12} live in an FBI file that was
lawfully released and is now hosted on the Bureau's own website, and they are forgeries; the Bureau
said so on the file itself. ``I have the FOIA documents'' and ``I have evidence'' are different
claims, and the gap between them has been the working capital of the paranormal publishing industry
for fifty years.

The Act's other limits are structural. It reaches federal agency records, so it does not reach
Congress, the courts, the President's own files, or --- critically --- a contractor's proprietary
records held on private premises. Special access programmes can be administratively walled so that
the office receiving the request genuinely has no responsive record to find, which is why a truthful
``no records'' letter proves far less than it appears to. Fees, backlogs measured in years, and the
releasing agency's own authority to decide what the exemptions cover do the rest. The requester is
asking the holder of the secret to adjudicate whether the secret is a secret.

Which is exactly why the events of Section~\ref{sec:121617} matter so much. For fifty years the
freedom-of-information request was the only instrument civilian researchers possessed --- a slow,
adversarial, one-page-at-a-time crowbar operated against an institution that controlled both the
records and the rules. What changed on 16 December 2017 was not that the crowbar got better. It was
that, for the first time, the government began generating the record and handing it over on its own
initiative, under statute, to a designated office with an obligation to report. The FOIA era is the
reason we know how much was withheld. The era that follows it is the first one in which somebody
inside the building is required to write the report anyway.

% ---------------------------------------------------------------
\section{Art Bell}\label{sec:artbell}

Arthur William Bell III was born on 17 June 1945 in Jacksonville, North Carolina, the son of a Marine
Corps captain and a Marine drill instructor, and he was a radio obsessive essentially from childhood
--- a licensed amateur operator as a boy, callsign eventually W6OBB. He served in the United States Air
Force during the Vietnam War as a medic, and he ran a pirate station from his barracks in Amarillo
before being stationed overseas.

And then there is the record. While working as a disc jockey at KSBK in Okinawa, Japan --- at the time
the only English-language commercial station in the country --- Art Bell set a Guinness world record by
broadcasting solo for 116 straight hours. One man, on air, alone, for very nearly five days.

I want to sit with that for a moment, because it is not a piece of trivia. It is the single most
explanatory fact about him. Nobody talks continuously for 116 hours unless the act of talking to
strangers over a transmitter is the thing they are actually for. Every characteristic of the show he
would later build --- the willingness to let a caller run, the tolerance for silence, the absence of a
screening panel and a producer's clock, the ability to hold five hours of unstructured live
conversation together night after night --- is visible in that record decades before \emph{Coast to
Coast AM} existed. The record was not a stunt he pulled. It was a description of him.

While in Okinawa he also, characteristically, involved himself in something with no professional
upside: he chartered aircraft and helped move Vietnamese orphans out of the country during the
war's collapse. He later worked in cable television management in Nevada, which is where the radio
career restarted.

What made him extraordinary as a broadcaster was a specific and rare skill, and it is one I have been
teaching in a different vocabulary for twenty chapters. Art Bell took guests seriously without
endorsing them. He would let a man explain, at length and without interruption or ridicule, that he had
been taken aboard a craft --- and he would ask precise follow-up questions, and he would not say he
believed it. He described himself as a sceptic who found the subject fascinating. That posture ---
genuine curiosity plus withheld assent --- is exactly the ``doubt-driven scepticism'' of Section~\ref{sec:doubtdrivenskepticism},
and Bell practised it live, unscripted, for hours, in front of millions, better than most people
manage in writing with time to revise.

He also had the instrument for it. The voice, the desert studio in Pahrump with the transmitter on his
own property, the sign-on from ``the high desert and the great American Southwest'', the deliberate
intimacy of talking to one person at three in the morning. Radio at that hour is a peculiar
technology: the listener is alone, usually driving or unable to sleep, and the voice is inside the car
with them. Bell understood that better than anyone who has done the job.

He retired, unretired, retired again, and retired again --- repeatedly, and for reasons that were
variously personal, family-related, and never fully explained --- which frustrated his audience for
twenty years. He died at his home in Pahrump on 13 April 2018, on Friday the 13th, which he would
have enjoyed enormously.

\subsection{Before the name: \emph{West Coast AM}}\label{sec:westcoastam}

The programme did not begin as \emph{Coast to Coast AM}, and the earlier name matters, because it tells
you what the show was before it became an institution. Bell's all-night slot started at KDWN in Las
Vegas --- a 50,000-watt clear-channel station broadcasting out of the Plaza Hotel downtown --- and it
was called \emph{West Coast AM}. The title was a literal description of the signal's reach at night: a
clear-channel AM carrier skipping across the western half of the continent after dark, to truckers,
shift workers, and insomniacs from Alberta to Baja. Bell himself later called KDWN ``where the magic
began.''\footnote{Bell interviewed by \emph{Radio Ink}, 20 July 2015, on the launch of
\emph{Midnight in the Desert}: \srclink{https://radioink.com/2015/07/20/art-bell-returns-to-radio-tonight/}.}
It was renamed \emph{Coast to Coast AM} once the show went out over a national satellite feed and the
old geographic modesty no longer fitted, a rebranding Bell credited to himself and to Alan Corbeth of
the syndicator.

\figwrap[r]{0.42}{Bell's house outside Pahrump, Nevada, with the antenna towers that put the
transmitter on his own property.}{https://www.flickr.com/photos/baggyrinkle/}{fig:artbellhome}

I should flag a discrepancy in my own text rather than paper over it. The dates given in
Section~\ref{sec:coastcoastam} --- a 1984 start and a 1992 national rebrand --- follow the account I
first learned and which RationalWiki still carries. The reference sources I have since checked put
\emph{West Coast AM} on KDWN from about 1978 and the rename to \emph{Coast to Coast AM} at 1988. I have
left the original passage as written and I am not going to quietly correct it, because the honest
position is that this book's own chronology of its subject is contested by a couple of years at the
front end, and the reader should know that even the datable facts about a man who broadcast for four
decades are not uniformly agreed. What is not in dispute is the sequence: local Las Vegas graveyard
shift first, national syndication second, and the name change in between.

The KDWN years were also not a smooth ascent. Bell was fired from that station twice, and the reason on
at least one occasion was precisely the material this book is about. When he put the retired Lockheed
engineer and UFO enthusiast John Lear on the air, the switchboard jammed so completely that the station
could not take ordinary business calls --- an early, purely mechanical demonstration that the audience
for this subject was very much larger than the format for it. In 1998 he bought a licence of his own and
founded \hi{KNYE 95.1 FM} in Pahrump, which is why he could sign on from the ``Kingdom of Nye'' and mean
it as an address rather than a joke: he owned a transmitter in Nye County, Nevada.

\subsection{The other programmes}\label{sec:bellothershows}

\emph{Coast to Coast AM} is the show everyone remembers, but Bell built and hosted several others, and
they are worth listing because each one tests a different part of the same format.

\hi{\emph{Area 2000}} came first, in 1993 --- a Sunday-evening programme underwritten by the Las Vegas
property developer Robert Bigelow, whose money recurs throughout this book and whose ranch data I
complain about in Section~\ref{sec:skinwalker}. It was explicitly a research-flavoured show, pointed at
UFO reports and funded by a man who wanted them investigated, and it is the direct ancestor of what
came next.

\hi{\emph{Dreamland}} replaced it as the Sunday-night slot and ran alongside the weeknight programme for
years, keeping the paranormal-and-conspiracy remit while the weeknight show broadened. Bell hosted it
and later handed it to Whitley Strieber, the \emph{Communion} author discussed earlier in this chapter
--- which is a neat illustration of the distribution network I keep describing, since the man whose book
fixed the modern image of the grey ended up holding the microphone.

\hi{\emph{Dark Matter}} was the failed comeback. Bell launched it on SiriusXM's Indie Talk channel on
16~September 2013 with the physicist Michio Kaku as his first guest, and it lasted six weeks; the last
programme aired on 4~November 2013. The dispute was about carriage and promotion rather than content,
and Bell said afterwards that the collapse cost him two years he would rather have spent on air. The
name survives in the \hi{Dark Matter Digital Network}, the streaming operation run by his long-time
engineer and webmaster Keith Rowland.

\hi{\emph{Midnight in the Desert}} was the last one, and the most interesting technically. It debuted on
20~July 2015 with Graham Hancock as the opening guest, carried by roughly 45 stations plus a TuneIn
stream and his own website, produced entirely out of Bell's own studio at the Pahrump house --- six
ZipLine phone circuits and two computers running Skype, which is the whole of the infrastructure a show
reaching a national audience actually needs once the satellite is optional. There was a paid tier at
five dollars called the Wormhole, a paranormal newscast at the bottom of the hour, and Crystal Gayle on
night one, whose song gave the programme its title. It ran five months. Bell walked away on
11~December 2015, citing the security of his home and family after a series of incidents there, and that
was the end of him as a working broadcaster.

\subsection{The lines}\label{sec:belllines}

The architecture of the show was its telephone system, and I mean that literally. Bell ran a bank of
named lines --- by his own account something on the order of fifteen circuits into the house, though I
have never found a published figure I would stake a citation on --- and each one was announced by name
rather than by number:

\begin{itemize}
  \item \hi{First Time Caller} --- reserved for people who had never been on the air before.
  \item \hi{East of the Rockies} --- the eastern half of the continent.
  \item \hi{West of the Rockies} --- the western half.
  \item \hi{International} --- everyone else, which on a clear-channel night meant a great deal of
        everyone else.
  \item \hi{Wildcard} --- no criteria at all.
\end{itemize}

That list looks like housekeeping and it is actually an editorial policy. A first-time-caller line is a
guarantee that the show cannot be captured by regulars. A wildcard line is an admission that the host
does not know in advance what is worth hearing. And for most of the programme's history there was no
screening at all --- callers went straight to air, unfiltered, until screening was finally introduced
around 2006 under a later host. From the standpoint of Chapter~\ref{ch:language} that is a machine for
generating unverifiable primary material at industrial scale, and from the standpoint of radio it is the
single bravest production decision in the medium. Both things are true. The scale is worth keeping in
view while you decide how you feel about it: 328 affiliate stations by February 1997, roughly 460 by
that June, and at peak more than 500 stations and something like fifteen million listeners a week.

\subsubsection{The themed lines}

The five circuits above were the permanent architecture. The part of the system people actually
quote came from the other half of the practice: on a themed night Bell would take a line off its
normal duty and reserve it for one category of caller, announce the category on air, and then let
whoever rang in have the microphone unscreened. The names were his, they were not fixed, and no
complete list exists, because they were invented for the night and retired with it. The recurring
ones, as far as the surviving audio and the show's own promotion support:

\begin{itemize}
  \item \hi{The Time Traveller line} --- open only to people claiming to have travelled in time,
        or to have met someone who had. This is the line that produced the John Titor material of
        Section~\ref{sec:johntitor} and a long tail of imitators.
  \item \hi{The Men in Black line} --- reserved for people claiming to have been visited,
        threatened or silenced after a sighting.
  \item \hi{The Werewolf line} --- for shapeshifter, dogman and Bigfoot-adjacent encounters, run
        on the nights given over to cryptids.
  \item \hi{The Ghost line}, which became the annual \emph{Ghost to Ghost AM} Halloween broadcast
        --- nothing but listeners' own hauntings, no guest, all night.
  \item \hi{The Weird or Strange People line} --- the catch-all for callers whose story fitted no
        category at all, which on this programme is a substantial constituency.
  \item \hi{The Whistleblower or Insider line} --- for people claiming direct knowledge from inside
        a government or contractor programme, sometimes with the caller's voice disguised and no
        name taken.
  \item \hi{The Open Lines} nights themselves, where every circuit was released at once and the
        show had no subject beyond whatever rang first.
\end{itemize}

I list these because they are the mechanism, not the decoration. \hi{A named line is a category,
and a category is an invitation: announce that the werewolf line is open and you will receive
werewolf calls, in volume, from an audience of millions, none of them screened and none of them
verifiable.} That is not an accusation of dishonesty against anybody involved. It is the plainest
demonstration in the whole of this book of the point Chapter~\ref{ch:language} makes about
collection method: what you ask for determines what arrives, and a body of testimony gathered this
way tells you a great deal about the audience and almost nothing about the phenomenon. The
counterweight is that the same open microphone is the only place where a great many of these
accounts were ever recorded at all, by anyone, and a field that had been abandoned by the
institutions in Chapter~\ref{ch:government} had nowhere else to put them.

\subsubsection{Two calls}

Two of those unscreened calls have outlived everything else the programme did, and they are the reason
the line names are quoted like scripture by people who have never heard the show.

The first arrived on the western line in 1997, and Bell's ``\emph{West of the Rockies}, you're on the
air!'' introduced a caller giving the name Mel Waters, who described a bottomless hole on his property
near Ellensburg, Washington --- nine feet across, walled with ancient brick, of no measurable depth,
with a dead dog thrown in that reportedly came back. \hi{Mel's Hole} has since been searched for by
geologists, journalists and county records offices, and there is no such property, no such hole, and no
Mel Waters in the assessor's rolls. It does not matter. The hole is now a fixture of American folklore,
endlessly retold, and it entered the record because a man phoned a radio station and nobody stopped him.

The second arrived on the eastern line in September 1997. Bell's ``\emph{East of the Rockies}, you're on
the air!'' brought a frantic, hyperventilating young man who said he was a former employee at Area 51,
that he had fled, that the things there were extradimensional beings, and that the government was about
to seize the emergency broadcast system --- and then, mid-sentence, the show went off the air. The
satellite uplink failed and stayed dead for the better part of twenty minutes. It is, purely as radio,
the most effective accident I have ever heard: a man warning of a conspiracy is silenced live by an
equipment fault, in front of millions, with no possible way for the audience to distinguish the two. The
caller later admitted the whole thing was a hoax, and Bell said the outage was a genuine technical
failure at the uplink. You should also know that this call is very widely misdated online to
28~April 1998; the September 1997 date is the one the reference sources support, and the confusion is
itself a nice small demonstration of how quickly a well-told story loses its metadata.

\vspace{6pt}
\pullquote[Art Bell, sign-on]{``From the high desert and the great American Southwest, I bid you all
good evening, good morning, good afternoon, wherever you may be in the world's time zones.''}

% ---------------------------------------------------------------
\section{Coast to Coast AM}\label{sec:coastcoastam}

Radio has proven to be one of the most successful and resilient technologies ever contrived.
Having survived the development of not one, but two advancements in technology that were otherwise
poised to render radio obsolete. First with the coming of television, and next with the inception
and later explosion of the internet nearing the end of the 20th century. Yet even today; more than
20 years later, radio is at the forefront of media consumption not drowned out by, but rather along
side of the internet and TV. Radio may even be attributed as the inspiration for the widely popular
spread of podcasts as a form of entertainment.

Podcasts have exploded in popularity recently, with many new additions showing up daily, trying to
emulate --- reproduce, if you will --- the form of other popular shows such as \emph{The Joe Rogan
Experience}, \emph{Freakonomics}, or \emph{Serial}. Podcasts have something cool, something unique
about them that separates and distinguishes them from the rest of radio. They exist to serve a niche,
fill a void, occupy some curiosity of the human condition that otherwise isn't served through
traditional channels. Podcasts are NOT --- well, some may be --- simply nothing but news, weather or
sports. They are a breath of fresh air, and for many the discovery that such a media exists is like an
escape from the monotonous day to day they have known their whole life.

Face it: news is boring. Weather and sports are not much more appealing. At least in the form of radio,
let me be clear. If you are going to listen to the announcer of your favourite sports match, would it
not be more entertaining, or just as easy, to pull up the stream on YouTube or through your TV and
receive visuals as well? If you are going to hear how Trump is single-handedly running the nation while
saving all the starving kids of Africa and finding a cure to cancer while he sleeps, then why not just
read that online? While things of this nature have always existed as a part of radio, and been consumed
through channels on the FM dial, AM radio is different. It's like a wildcard, a stab into the darkness,
and an alternative world in the space of NPR --- National Public Radio. And in the more than
hundred-year history of radio stations broadcasting content into the great unknown, none other has
been as provocative, as influential, and downright weird as \emph{Coast to Coast AM}.

Beginning in 1984, at a time when nearly all non-music orientated radio stations were focused on
covering the standards --- that being news, politics, and sports --- \emph{Coast to Coast} was
different. It began in the desert of Las Vegas when a man by the name of Art Bell was approached with
an opportunity to move his daytime talk show to the wee hours of the night. Thrilled with the chance,
Art gleefully accepted, and while currently working as the manager of a nearby cable company, he
quickly found himself pulled back into the world of radio he grew up on. While initially running as
another typical conservative-orientated talk show, Art would ever so often inject things of his own
curiosity into the program. This proved extremely successful, in addition to improving the ratings of
the show, and so in 1992 Art Bell's all-night show was rebranded as \emph{Coast to Coast AM} and taken
nationwide.\footnote{See \srclink{https://rationalwiki.org/wiki/Coast_to_Coast_AM}.} While \emph{Coast to
Coast} was in fact an FM broadcast station, it mirrored much of the content style of things found on
the AM dial. Instead of music, it was people talking for hours upon hours on whatever the subject at
hand may be, all under the guidance of the show's host.

The subject of choice you ask? Well anything and everything from the occult, to ghosts and demons,
aliens and UFO's, crop circles, and the Bermuda triangle. Ley lines and ancient history, people who
claim to be abducted, people who claim to be abducting, vampires, werewolves, men who superglue
their mouth shut, the men in black, people pursued by the men in black, people pursuing the men in
black, some guy who collects toenails, Roswell \& alien autopsies, government secrets and Project
Blue Book, space, time, the whole nine yards and yes everything INCLUDING the kitchen sink!!! In
other words that show was WEEEIRD!

\begin{funfact}[Fun Fact: The Host Glued His Own Mouth Shut]
That list is not padded, and I want to defend one item on it in particular, because the man
who superglued his mouth shut was Art Bell. During a break he was repairing a shelf in the
studio, got superglue on his fingers, and in the process of dealing with it managed to
seal his own lips together --- and then had to come back from the break and explain, live on
air, to several million people, that the host of the biggest overnight talk show in America
could not currently talk. He turned it into a station promo, which is the only surviving
artefact of the episode I can point you at: an aircheck dated 21 February 1997 in which Bell
announces, with the wounded dignity the situation deserves, that he has glued his lips
together.\footnote{Coast to Coast AM promo, 21 February 1997:
\srclink{https://www.podomatic.com/podcasts/frankomundo/episodes/2013-11-25T16_16_47-08_00}.}
I include it because it is a perfect miniature of why the programme worked. A conventional
broadcaster buries that. Bell put it on the air.
\end{funfact}

What is worth adding, from the perspective of the twenty chapters behind us, is how much of this book
would not exist without that show. Chapter~\ref{ch:earthfiles} is largely about Linda Moulton Howe,
whose reporting reached a mass audience primarily as a recurring \emph{Coast to Coast} guest.
Section~\ref{sec:dulce} and the underground-base literature, the Mothman revival, remote viewing, the
Hale-Bopp companion object of 1997 --- which Bell later publicly acknowledged handling badly, and
which intersected tragically with the Heaven's Gate deaths discussed in
Chapter~\ref{ch:religion} --- and a substantial fraction of the named figures in this text either
broke their material on that programme or became known through it. Art Bell did not merely report on
the field. Between 1992 and 2003 he was very nearly the field's distribution network.

\subsection{George Noory}

\figwrap[r]{0.34}{George Noory at Alien Snowfest in Big Bear, California, 2019. He has hosted \emph{Coast to Coast AM} since 2003, which is longer than Bell hosted it.}{}{fig:noory}


Bell retired --- for what turned out to be one of several times --- at the end of 2002, and the chair
went to George Noory. It is worth knowing how he got there, because the story is a small case study in
how this subject actually propagates.

Noory was born in Detroit in 1950 and came up through television news, not through paranormal radio. He
started as a newscaster at WCAR-AM, produced news at WJBK-TV in Detroit through the mid-1970s, and went
on to run newsrooms as news director at KMSP-TV in Minneapolis and KSDK-TV in St.\ Louis, collecting
three local Emmy Awards along the way. He also served nine years as a lieutenant in the United States
Naval Reserve. That is a conventional broadcast-journalism résumé, and it matters: the man who has
spent the last two decades as the field's largest single microphone was trained to write copy for a
six-o'clock newscast.

The unconventional half was there from the beginning. Noory has said he became interested in the
paranormal and in UFOs as a child, and as a teenager he joined NICAP --- the National Investigations
Committee on Aerial Phenomena, the largest civilian investigative body of its era. In 1996 he
began hosting a late-night open-phones programme called \emph{The Nighthawk} on KTRS in St.\ Louis. It
was that show which caught the attention of Premiere Radio Networks, the syndicator of \emph{Coast to
Coast AM}. He was brought in as a guest host in April 2001, took over the Sunday-night edition from
Ian Punnett, and in January 2003, on Bell's retirement, inherited the weeknight programme
outright.\footnote{See \srclink{https://en.wikipedia.org/wiki/George_Noory}.} He has held it ever since,
broadcasting out of Los Angeles, which is why Bell's sign-off placed the show simultaneously ``from the
kingdom of Nye and the city of angels.'' Outside the show he has co-written several books, appeared
recurrently on the History Channel's \emph{Ancient Aliens}, and fronted the streaming series
\emph{Beyond Belief}.

The critical assessment of Noory, and I include it because this is a book about how claims survive
transmission, is that he is a gentler filter than Bell was. \emph{The New York Times} put it plainly:
Bell, if he thought a caller was ``loo-la,'' would say so, and Noory generally does not. Whitley
Strieber, a frequent guest, framed the same trait as a virtue --- Noory is ``willing to take the
intellectual journey with a guest'' and let the audience find the kernel worth keeping. Both
descriptions can be true at once, and the difference is not really about the two men. It is about what
an open microphone does to the record. A host who pushes back leaves a transcript in which the weak
claim is visibly weak; a host who does not leaves a transcript in which it is merely a claim, and
claims travel. That is the same mechanism as the ``Ghost to Ghost'' problem in the box below, with a
longer fuse.

What is not in dispute is the scale. Noory has now hosted the programme for well over twenty years ---
substantially longer than Bell hosted it nationally --- across hundreds of affiliate stations. For a
whole generation of listeners, including a large share of the students I have taught, \emph{Coast to
Coast AM} means Noory's voice and not Bell's. And I have a personal reason for saying so, which the
reader has already met in the preface: he had me on the air in the autumn of 2017 to talk about what a
Canadian university was doing teaching this subject at all, and it was Noory who told me, a few months
later and again live on air while I sat in a University of Lethbridge computer lab drawing the cover of
this book, that Art Bell had died.\footnote{The author's live conversation with George Noory on \emph{Coast to Coast AM}:
\srclink{https://www.dropbox.com/scl/fi/g3kgjxuoejislpwtq6bbm/U-of-L-Ufology-on-Coast-to-Coast-AM.mov?rlkey=r8apwt11bksvihbrxl4en4a99}.}
I am, in that narrow sense, part of the distribution network I have just spent a chapter describing,
and I would rather declare it than pretend to a vantage point outside the thing I am studying.

\begin{funfact}
Bell's open-phones formats are an accidental monument to everything in Chapter~\ref{ch:language}.
``Ghost to Ghost AM'' on Halloween took unscreened listener stories on air. There was no verification,
no follow-up, no possibility of investigation --- and that was the point; it was theatre, and Bell said
so. But those broadcasts entered the literature anyway, cited later by authors who no longer
remembered that the origin was an unscreened Halloween call-in. That is how anecdote becomes source
material: not by anyone lying, but by the label falling off.
\end{funfact}

\subsubsection{Digital Book Burning: The Art Bell Vault}

I have one grievance with this programme and it is not about credulity, so let me set it out plainly,
because it is the only place in this book where I am angry at an institution I otherwise admire.

What Art Bell built between 1992 and 2003 was not a product. It was a commons. The whole architecture
of the show --- open lines, no screening on some nights, first-time callers, the wildcard line, a
truck driver in Nebraska given the same four minutes as a professor --- was a machine for letting
anybody speak into the American night and be heard by millions. That is the thing that made it
historically important, and it is the reason it appears in a textbook at all. Twenty chapters of this
book trace claims that entered the record through those microphones. The tapes are not entertainment
assets. They are the primary source material for a generation of belief, and they are the closest
thing this field has to a spoken archive of its own formation.

They are now behind a paywall. In September 2019 \emph{Coast to Coast AM} launched the
\hi{Art Bell Vault} as a benefit for \term{Coast Insider} subscribers --- a curated, commercial-free
run of vintage programmes, opening with twenty shows and adding two a week, available only to people
paying roughly \$6.95 a month or \$54.95 a year through the Coast app and website. Alongside that, the
free copies went away. For years the classic broadcasts sat on YouTube, on Spotify, on fan-run
archives, uploaded by listeners who taped them off the air at two in the morning and who were, in the
most literal sense, preserving them. I listened to hundreds of hours of Art Bell that way, for nothing,
as a student, and a good deal of my education in this subject came out of those uploads. Then the
takedowns began, and channel after channel and playlist after playlist was systematically purged.

I call that \hi{digital book burning}, and I am aware that the phrase is strong. I mean it as a
description rather than an insult. Nobody destroyed a master tape. But the practical effect of pulling
every free copy out of circulation while metering the originals two shows a week to paying subscribers
is that a body of knowledge which was in the public's hands is no longer in the public's hands. A
curated vault is not an archive. An archive is complete, browsable, and citable; a vault is a
release schedule, and somebody else decides which episodes are worth releasing and which are quietly
never selected. A scholar cannot cite a subscription tier. A listener in a country the app does not
serve, or a listener who does not have seven dollars a month, has simply been removed from the
audience of a show whose entire founding premise was that everyone gets a turn at the microphone.

In fairness to the man, the copyrights sit with Premiere Radio Networks and its parent, not with the
host, and a syndicator monetising its own library is doing nothing a court would object to. But there
is a custodianship here that is separate from ownership, and it belongs to whoever holds the chair. My
reader should understand where the keys to this castle actually hang. Not in Washington --- not in an
Air Force records office, not behind a classification stamp, not in any of the vaults I spend the rest
of this book complaining about. The largest single body of primary-source material in modern ufology is
locked up by the show that made it, and the man with the keys is George Noory.

\begin{funfact}
\funfactlead{Notice the shape of that complaint.} For twenty chapters I have argued that the field's
deepest injury is not the absence of evidence but the enclosure of it: material that exists, that
somebody holds, that you are not permitted to examine. I have made that case against the Air Force,
against contractors, against Bigelow's ranch data in Section~\ref{sec:skinwalker}, and against every
researcher who promises documents next year. The uncomfortable finding is that ufology does it to
itself as readily as any government does it to ufology --- and that when the gate goes up, the
reason is almost never conspiracy. It is a subscription business model. Keep that in mind the next
time somebody tells you the truth is being suppressed. Sometimes it is merely being sold.
\end{funfact}

% ---------------------------------------------------------------
\section{12-16-17}\label{sec:121617}

On 16 December 2017, \emph{The New York Times} published a story titled ``Glowing Auras and Black
Money: The Pentagon's Mysterious U.F.O.\ Program,'' reporting the existence of the Advanced Aerospace
Threat Identification Program, a Defense Department effort funded at roughly \$22 million between 2007
and 2012, largely at the direction of Senator Harry Reid. Alongside it the paper published gun-camera
footage from Navy aircraft, including the \emph{Nimitz} encounter described in Section~\ref{sec:uap}.

\figfull[3.1in]{The consequence of that morning, six years on: David Grusch (left), Ryan Graves
and David Fravor sworn in before the House Oversight subcommittee on 26~July 2023. Before
16~December 2017 no such hearing had been held in over fifty years, and no serving officer had
a statutory channel through which to testify at all. Whatever one makes of the claims, the
room itself is the change.}{}{fig:grusch}

I have given this date its own section, and the date is the section title, because it is the dividing
line. Everything in this book belongs to one side of it or the other.

Before 16 December 2017, the structure of the subject was fixed and had been for seventy years. Civilian
researchers gathered testimony; the government's public position was that there was nothing to
investigate and had not been since Blue Book closed in 1969; the burden fell entirely on believers to
produce evidence; and the only available channels were paperbacks, conferences, late-night radio and
the internet. Chapters~\ref{ch:government} through \ref{ch:nuclear} describe that world.

After that date, the Department of Defense had confirmed in a newspaper of record that it had run a
programme, that it possessed sensor data it could not explain, and that it considered the question a
matter of flight safety and airspace security. Within four years there was a statutory reporting
office, an unclassified Office of the Director of National Intelligence preliminary assessment in June
2021, congressional hearings for the first time in over fifty years, and the term \term{UAP} in the
United States Code.

What actually changed was jurisdiction. The question moved out of the hands of enthusiasts and into the
hands of institutions with sensors, budgets, subpoena power and an obligation to report. That is a
smaller and more bureaucratic transformation than the excitement of that week suggested --- and it is
also a far more consequential one, because for the first time the subject acquired the one thing
Chapter~\ref{ch:language} said it never had: an official definition, a designated custodian, and a
requirement to produce data rather than testimony.

So the old guard and the new guard differ less in what they believe than in what they consider
evidence. The old guard was built on witnesses, because witnesses were all anyone had; its craft was
interviewing, archiving, and the freedom-of-information request, and its heroes were researchers who
spent decades on cases the state refused to acknowledge. Its besetting weakness was that testimony
accumulates without ever converging, which is why eighty years produced a vast literature and no
settled fact. The new guard is built on instrumentation --- radar tracks, infrared imagery, calibrated
sensor fusion --- and its besetting weakness is the opposite: sensors have artefacts, and
Section~\ref{sec:uap} showed how much of the celebrated Navy footage is arguably explicable as
parallax, glare, and gimbal rotation. Neither guard has the answer. But only one of them has
instruments capable of producing one, and that is why the shift matters.

\subsection{Immaculate Constellation}\label{sec:immaculate}

There is a coda to 16 December 2017 that only surfaced seven years later, and it belongs here rather
than in Chapter~\ref{ch:government} because the allegation is not merely that a secret programme
exists --- it is that the programme was created \emph{in reaction to that morning's newspaper}.

On 8 October 2024, the journalist Michael Shellenberger published on his newsletter \emph{Public} an
eleven-page report he had received from a source he described as a current or former United States
government official acting as a whistleblower. The report alleged the existence of an unacknowledged
special access programme --- a \term{USAP}, meaning a programme whose very existence is classified ---
codenamed \textsc{Immaculate Constellation}. According to the document, the Department of Defense stood
the programme up in 2017, after the \emph{New York Times} AATIP story, to consolidate UAP imagery and
observations gathered by tasked and untasked collection platforms --- that is, by satellites, aircraft
and sensors both deliberately pointed at the phenomena and incidentally capturing them --- and that it
has operated outside congressional oversight ever since. The report organised its material into seven
categories of evidence and cited specific encounters, among them an F-22 intercepted by orbs and Navy
personnel observing an orange-red sphere. Shellenberger testified that a source had warned him that
merely printing the codename could expose the publisher to surveillance under the Foreign Intelligence
Surveillance Act, and that the secrecy was enforced ``with a lot of vigor.''

On 13 November 2024, the House Oversight Committee's subcommittee on National Security held a hearing
titled \emph{Unidentified Anomalous Phenomena: Exposing the Truth}. Shellenberger appeared alongside
Luis Elizondo, the former head of AATIP, and Rear Admiral Tim Gallaudet, USN (ret.), who testified that
in January 2015 he had received an email on a secure Navy network headed ``URGENT SAFETY OF FLIGHT
ISSUE'' about near-midair collisions during a \emph{Theodore Roosevelt} strike group exercise, with a
now-declassified video attached --- and that the email vanished from his inbox the next day and the
matter was never raised again. Representative Nancy Mace entered the \textsc{Immaculate Constellation}
document into the Congressional Record, which is how an unsigned eleven-page PDF became a citable
artefact. In May 2025 a former Pentagon metadata analyst and State Department intelligence liaison named
Matthew Brown publicly identified himself as its author.

The Pentagon's denial is unusually specific, and specificity is worth noticing in a field where official
language is normally built to leak. A Department of Defense spokeswoman, Sue Gough, stated that the
Department ``has no record, present or historical, of any type of SAP called `IMMACULATE
CONSTELLATION','' that all classified UAP programmes are properly briefed to Congress, and that the
circulating document is not an official DoD or ODNI product. AARO has separately attributed the broader
family of secret-programme claims to \emph{circular reporting} --- a closed group of people who each
believe the rumour partly because the others do, so that a single origin is mistaken for independent
confirmation.

I have no way to adjudicate this and neither does anyone else outside a cleared reading room, so let me
say plainly what the situation actually is. What exists in public is a document with no signature at the
time of release, describing imagery that has not been released, sourced to interviews and recollection
rather than to the imagery itself. \hi{That is testimony about testimony again, which is the oldest
structure in this book,} and the rule from Section~\ref{sec:correlationvscausation} does not get suspended because the codename is a
good one. It \emph{is} a good one --- \textsc{Immaculate Constellation} sounds like something that must
be real, and a name that does that much rhetorical work on its own deserves more suspicion than less.
Equally, the denial is not proof either: a USAP is by construction a thing officials are obliged to deny,
which is precisely the epistemic trap Section~\ref{sec:grusch} described. If the claim is true, classification is why
we cannot check it; if it is false, classification is also why we cannot check it.

What makes this the right place to end the section is that \textsc{Immaculate Constellation} is a test
case for the thesis of the chapter rather than a footnote to it. The 2017 story moved the subject because
somebody produced a record and filed it through a channel that could not ignore it. The 2024 allegation
has not moved it, despite a hearing, a Congressional Record entry and a named author, because what was
filed was a description of records rather than the records. Same subject, same committee, same city ---
different evidentiary object, and therefore a different outcome. Watch which one gets you a statutory
office and which one gets you a denial and a podcast tour.

A personal note, since this is the last chapter and I have earned it. This happened just a few months
after I started the U of L Ufology Group. I had spent that autumn lecturing to a room of six people ---
down from the hundred-odd who had signed the mailing list --- on a subject I was regularly told was
beneath serious discussion. Once a semester we got out of that room and did the thing the subject is
actually about: a field trip to the Old Man Observatory, south of Lethbridge, where the Ufology Group
and the Lethbridge Astronomy Society stood in the cold together and looked at the sky through a
telescope instead of arguing about it indoors.\footnote{The Old Man River Observatory, operated by the
Lethbridge Astronomy Society: \srclink{https://lethbridgeastronomysociety.ca/the-observatory.html}.}
It is a detail worth keeping, because the two groups are often assumed to be enemies and were not:
the astronomers knew what was in the sky, we wanted to know what was not accounted for, and one
evening at an eyepiece teaches more about the difference than a term of lectures. And then in the middle of December the Pentagon confirmed it had been
quietly funding exactly the thing that did not warrant discussing. I would like to report that I
handled this gracefully. I did not; I was insufferable for approximately a month. What I did not
understand at the
time, and what this whole book has been my attempt to work out, is that the story did not vindicate
believers at all. It vindicated \emph{procedure}. The reason we now have hearings and a reporting
office is not that someone finally believed the witnesses. It is that somebody produced a sensor
record, wrote it up, and filed it through a channel that could not simply ignore it.

\hi{That is the whole thesis of this textbook in one sentence.}

\actualq{whether something is out there.}{whether we are willing to build the tools, the
definitions and the discipline required to find out --- and to accept the answer if it turns
out to be boring.}

Eighty years of testimony moved this subject almost nowhere. One reporting requirement moved it further
in four years than the previous seventy combined. Keep looking up. Then write down what your instrument
said, not what you hoped it said.

\vspace{14pt}
\pullquote[Art Bell, sign-off]{``And so, from the high desert, good night.''}

\begin{srcnotes}
\item Freedom of Information Act, 5 U.S.C.\ \S~552 --- signed 4 July 1966, effective 5 July 1967; the
nine exemptions, including b(1) and b(3).
\item Privacy Act and Freedom of Information Act Amendments of 1974, Pub.\ L.\ 93--502 --- enacted over
President Ford's veto; response deadlines, fee limits, attorney's fees, and \emph{de novo} judicial
review with authority for \emph{in camera} inspection of withheld records.
\item \emph{Ground Saucer Watch, Inc.\ v.\ Central Intelligence Agency} --- litigation begun 1977;
approximately 900 pages released December 1978, 57 documents withheld under exemptions b(1) and b(3).
\item \emph{C.I.A.\ Papers Detail U.F.O.\ Surveillance}. \emph{The New York Times}, 14 January 1979 ---
contemporaneous reporting of the Ground Saucer Watch release.
\item \emph{Citizens Against UFO Secrecy v.\ National Security Agency}, D.D.C.\ 1980 --- 156 documents
identified, 135 withheld; \emph{in camera} affidavit of Eugene F.\ Yeates, Chief of Policy, NSA,
filed 9 October 1980; redacted affidavit publicly released 2005.
\item \emph{Phillippi v.\ Central Intelligence Agency}, 546 F.2d 1009 (D.C.\ Cir.\ 1976) --- origin of
the ``neither confirm nor deny'' (Glomar) response.
\item Federal Bureau of Investigation, \emph{The Vault} --- ``Majestic 12'' file, released under FOIA
and the source of Figure~\ref{fig:mj12memo}. \srclink{https://vault.fbi.gov/UFO}
\item Cooper, H., Blumenthal, R.\ \& Kean, L.\ (16 December 2017). Glowing Auras and Black Money: The
Pentagon's Mysterious U.F.O.\ Program. \emph{The New York Times}.
\item Office of the Director of National Intelligence (25 June 2021). \emph{Preliminary Assessment:
Unidentified Aerial Phenomena}.
\item Shellenberger, M.\ (8 October 2024). Pentagon Is Illegally Hiding Secret UFO Program From
Congress, Say Whistleblowers. \emph{Public}. \srclink{https://www.public.news/p/pentagon-is-illegally-hiding-secret}
\item U.S.\ House of Representatives, Committee on Oversight and Accountability, Subcommittee on
National Security, the Border, and Foreign Affairs (13 November 2024). \emph{Unidentified Anomalous
Phenomena: Exposing the Truth}. Hearing transcript and witness statements (Elizondo, Gallaudet,
Shellenberger); includes the eleven-page \textsc{Immaculate Constellation} document entered into the
record by Rep.\ Nancy Mace.
\item Department of Defense, statement of spokeswoman Sue Gough (October 2024), reported
contemporaneously; the Department has ``no record, present or historical'' of a special access programme
of that name.
\item \emph{Coast to Coast AM}. RationalWiki. \srclink{https://rationalwiki.org/wiki/Coast_to_Coast_AM}
\item Premiere Radio Networks (September 2019). ``\emph{Coast to Coast AM} Launches the Art Bell Vault
for Coast Insiders'' --- press release announcing the paid vault, its opening twenty commercial-free
vintage programmes and the two-per-week release schedule.
\srclink{https://www.premierenetworks.com/press/coast-coast-am-launches-art-bell-vault-coast-insiders}
\item \emph{Coast to Coast AM}, Coast Insider subscription pages and billing help centre, for the
current monthly and annual subscription rates quoted in Section~\ref{sec:coastcoastam}.
\srclink{https://www.coasttocoastam.com/pages/coastinsider/}
\item \emph{Art Bell}. Wikipedia --- for \emph{West Coast AM} on KDWN Las Vegas, the 1988 rename to
\emph{Coast to Coast AM} with Alan Corbeth, the move of the studio from the Plaza Hotel to the Pahrump
house, the founding of KNYE 95.1~FM in 1998, affiliate and audience figures, and the death at Pahrump on
13 April 2018. \srclink{https://en.wikipedia.org/wiki/Art_Bell}
\item \emph{Radio Ink} interview with Art Bell, 20 July 2015 (also carried by Puget Sound Radio), for
KDWN as ``where the magic began,'' the John Lear broadcast that jammed the station switchboard, the two
dismissals from KDWN, the SiriusXM \emph{Dark Matter} dispute, and the technical set-up of
\emph{Midnight in the Desert} --- six ZipLine circuits, two Skype computers, the \$5 Wormhole tier and
the bottom-of-the-hour paranormal newscast.
\srclink{https://radioink.com/2015/07/20/art-bell-returns-to-radio-tonight/}
\item \emph{Midnight in the Desert} and \emph{Dark Matter} programme entries, for the 16 September 2013
debut with Michio Kaku, the 4 November 2013 finale, the 20 July 2015 launch with Graham Hancock across
roughly 45 affiliates, and the retirement announcement of 11 December 2015 on home-security grounds.
\srclink{https://en.wikipedia.org/wiki/Midnight_in_the_Desert}
\item \emph{Mel's Hole}, for the 1997 ``West of the Rockies'' call and the subsequent failure of
geologists, reporters and county records to locate either the hole or the caller.
\srclink{https://en.wikipedia.org/wiki/Mel\%27s_Hole}
\item The ``Area 51 caller'' broadcast of September 1997 and the satellite-uplink failure that followed
it; the caller's later admission of a hoax, and the widespread misdating of the episode to 28 April 1998.
\item Washington Post (February 1997) and \emph{The Oregonian} (June 1997) profiles of Bell, for the 328
and 460 affiliate counts quoted in Section~\ref{sec:belllines}.
\item Photograph: Bell's house and antenna towers outside Pahrump, Nevada, by Flickr user
``baggyrinkle,'' 19 March 2006, reused under the Creative Commons licence stated on the original posting.
\item \emph{James Fox (film-maker)}. Wikipedia --- for the birth date of 25~April 1968, the family
account that brought him to the subject, and the filmography and dates used in the opening pages of
this chapter. \srclink{https://en.wikipedia.org/wiki/James_Fox_(filmmaker)}
\item \emph{Out of the Blue} (2003), dir.\ James Fox; narrated by Peter Coyote; Sci~Fi Channel
broadcast premiere 24~June 2003. Festival and theatrical material from 2002 accounts for the
widespread ``2002'' dating discussed in the Fun Fact.
\srclink{https://www.imdb.com/title/tt0357324/}
\item \emph{UFOs: 50 Years of Denial?} (1997), \emph{I Know What I Saw} (2009), \emph{Pretty Slick}
(2016), \emph{The Phenomenon} (2020, with Marc Barasch), \emph{Moment of Contact} (2022) and
\emph{The Program} (2024) --- production and release details from the films' own credits and
distributor listings.
\item Witness roster of \emph{Out of the Blue} and \emph{I Know What I Saw}, including Gordon Cooper,
former Belgian Air Force Chief of Staff Wilfried De~Brouwer, Robert Salas on the Malmstrom Echo Flight
shutdown, and Parviz Jafari on the 1976 Tehran incident; cross-referenced against the primary sources
cited in Chapter~\ref{ch:nuclear}.
\item Photograph: James C.\ Fox, by Commons contributor OldDominionVA89, CC~BY-SA~4.0.
\item Photograph: the ``ET de Varginha'' memorial, Varginha, Minas Gerais, February 2025, by
Commons contributor ``Ent\~ao Vamo!,'' CC~BY-SA~4.0. No photograph of the 1996 object or entity
exists; none is claimed here.
\end{srcnotes}
